
write the equation of the line satisfying the given set of conditions in slope-intercept form.

\begin{exercise}\label{exercise-2.2.1}
passing through $(2,6)$ and $(1,8)$
\end{exercise}

\begin{solution}
    $y=-2x+10$
\end{solution}



\begin{exercise}
passing through $(-4,7)$ and $(-10,-8)$
\end{exercise}

\begin{solution}
    $y=\frac{5}{2}x+17$
\end{solution}



\begin{exercise}
with slope $2.1$ passing through $(0.1,0)$
\end{exercise}

\begin{solution}
    $y=2.1x-0.21$
\end{solution}



\begin{exercise}
with slope $6$ passing through $(-5,2)$
\end{exercise}

\begin{solution}
    $y=6x+32$
\end{solution}



\begin{exercise}
with slope $-\frac{1}{2}$ passing through $(\frac{2}{5},-1)$
\end{exercise}

\begin{solution}
    $y=-\frac{1}{2}x-\frac{4}{5}$
\end{solution}



\begin{exercise}
passing through $(1,3)$ and $(-7,3)$
\end{exercise}

\begin{solution}
$y=3$
\end{solution}



\begin{exercise}
passing through $(5,-5)$ and $(5,2)$
\end{exercise}

\begin{solution}
    $x=5$
\end{solution}



\begin{exercise}\label{exercise-2.2.8}
with slope 0 and passing through $(-3,-1)$
\end{exercise}

\begin{solution}
    $y=-1$
\end{solution}





 In \exrange{exercise-2.2.9}{exercise-2.2.16}, write the formula in slope-intercept form for the linear function $f(x)$ satisfying the given set of conditions.

\begin{exercise}\label{exercise-2.2.9}
$f(2)=-5$ and $f(4)=0$
\end{exercise}

\begin{solution}
    $f(x)=\frac{5}{2}x-10$
\end{solution}



\begin{exercise}
$f(1)=-2$ and $f(3)=-5$
\end{exercise}

\begin{solution}
    $f(x)=-\frac{3}{2}x-\frac{1}{2}$
\end{solution}



\begin{exercise}
$f(x)$ has slope $2.1$ and its graph passes through $(0.1,0)$
\end{exercise}

\begin{solution}
    $f(x)=2.1x-0.21$
\end{solution}



\begin{exercise}
$f(x)$ has slope $-1$ and its graph passes through $(-1,2)$
\end{exercise}

\begin{solution}
    $f(x)=-x+1$
\end{solution}



\begin{exercise}
the graph of $f(x)$ passes through $(-1,3)$ and $(-7,-2)$
\end{exercise}

\begin{solution}
    $f(x)=\frac{5}{6}x+\frac{23}{6}$
\end{solution}



\begin{exercise}
the slope of $f(x)$ is 3 and its horizontal intercept is $-2$
\end{exercise}

\begin{solution}
    $f(x)=3x+6$
\end{solution}



\begin{exercise}
the slope of $f(x)$ is 1 and its vertical intercept is $2$
\end{exercise}

\begin{solution}
    $f(x)=x+2$
\end{solution}



\begin{exercise}\label{exercise-2.2.16}
the vertical intercept of $f(x)$ is $3$ and its horizontal intercept is $7$
\end{exercise}

\begin{solution}
    $f(x)=-\frac{3}{7}x+3$
\end{solution}






\begin{exercise}
Which of the following lines is NOT the graph of a linear function?
    \begin{parts}
        \item $\displaystyle y=-2x-3x+1$
        \item $\displaystyle y=-5$
        \item $x=7$
    \end{parts}
\end{exercise}


\begin{solution}
(c) is not the graph of a linear function.
\end{solution}





\begin{exercise}
The distance in miles from the finish line, $D(t)$, of a bicyclist $t$ hours after beginning a race is given by $D(t)=60-20t$
    \begin{parts}
        \item Write a complete sentence explaining the practical meaning of the vertical intercept of this linear function. Include units in your answer.

        \item Write a complete sentence explaining the practical meaning of the horizontal intercept of this linear function. Include units in your answer.

        \item Write a complete sentence explaining the practical meaning of the slope of this linear function. Include units in your answer.
    \end{parts}
\end{exercise}

\begin{solution}
\begin{parts}
    \item The vertical intercept 60 is the distance in miles that the biker is from the finish line at the start of the race.

    \item The horizontal intercept 3 is the number of hours that it takes the biker to complete the race.

    \item The slope $-20$ indicates that the distance between the biker and the finish line is decreasing at a rate of 20 miles per hour.
\end{parts}
\end{solution}



\begin{exercise}
The value of an antique lamp, $V(t)$, in dollars, $t$ years after its purchase is given by
\[
V(t)=1900+250t.
\]

\begin{parts}
\item What was the purchase price of the lamp?

\item When will the value of the lamp reach \$4000?
\end{parts}
\end{exercise}


\begin{solution}
\begin{parts}
    \item \$1900
    \item The lamp will have a value of 4000 dollars 8.4 years after its purchase.
\end{parts}
\end{solution}



\begin{exercise}
Two mobile phone companies sell an international roaming plan. Company A charges a $\$75$ dollar fixed monthly fee and $\$0.20$ per minute for talk. Company B charges a $\$90$ dollar fixed monthly fee and $\$0.10$ per minute for talk. 

\begin{parts}
\item Write the linear functions $A(m)$ and $B(m)$ which give the monthly cost charged by Company A and Company B, respectively, with $m$ minutes spent talking on the phone.

\item For what value of $m$ is the cost the same with either company?

\item Which company gives you a better deal if you plan to talk for hours?
\end{parts}

\end{exercise}


\begin{solution}
    \begin{parts}
        \item $A(m)=0.20m+75$; $B(m)=0.10m+90$
        \item $m=150$
        \item $B(m)$
    \end{parts}
\end{solution}







 In \exrange{exercise-2.2.21}{exercise-2.2.24}, solve each linear equation. If there is no solution or infinitely many solutions, say so.

\begin{exercise}\label{exercise-2.2.21}
$\displaystyle 0.1x-(2x-5)=2.5$
\end{exercise}

\begin{solution}
$x\approx1.3158$
\end{solution}




\begin{exercise}
$\displaystyle \frac{-2(x-3)}{5}=x+1$
\end{exercise}

\begin{solution}
$x=\frac{1}{7}$
\end{solution}



\begin{exercise}
$\displaystyle 3x-7=5+3x$
\end{exercise}

\begin{solution}
no solution
\end{solution}



\begin{exercise}\label{exercise-2.2.24}
$\displaystyle 2x+2=4+2(x-1)$
\end{exercise}

\begin{solution}
infinite number of solutions
\end{solution}